\documentclass[11pt]{article}
\usepackage[utf8]{inputenc}
\usepackage[T1]{fontenc}
\usepackage[margin=1in]{geometry}
\usepackage{hyperref}
\usepackage{enumitem}
\hypersetup{colorlinks=true,linkcolor=black,urlcolor=blue,citecolor=blue}
\title{Block Coding with Sphero}
\author{Piter Garcia}
\date{March 20, 2026}
\begin{document}
\maketitle
\section*{Packet Use}
This printable lesson packet is generated from the paired Markdown guide. The Markdown version remains the clickable, neurodivergent-friendly working copy; this PDF carries the same lesson substance in a formal handout format. Evidence claims in this packet are based on transcript-derived lesson notes and the cleaned transcript files in this workspace.

\section*{Quick Scan}
\begin{itemize}[leftmargin=*]
\item \textbf{Grade:} Grade 3
\item \textbf{Grade Hub:} Notes [../grade-3-Notes.md]
\item \textbf{Schedule Reference:} Spring2026\_TeachingSchedule\_Derek.md [../../school/Spring2026\_TeachingSchedule\_Derek.md]
\item \textbf{Workbook Sheet:} \texttt{3rd} in \texttt{Teaching\_Placement-IGNITE Curriculum\_ GCSD25-26.xlsx}
\item \textbf{Lesson Focus:} Students pair robots, create movement programs, use loops for repeated motion, and debug route accuracy through floor-based challenges.
\item \textbf{CS Alignment Strength:} Strong CS lesson
\item \textbf{LaTeX Source:} block-coding-with-sphero.tex [./block-coding-with-sphero.tex]
\item \textbf{Bib File:} block-coding-with-sphero.bib [./block-coding-with-sphero.bib]
\item \textbf{Artifacts Folder:} artifacts/ [./artifacts/]
\end{itemize}

\section*{Lesson Identity}
\begin{itemize}[leftmargin=*]
\item \textbf{Likely Class Blocks:} Day 3 -- 9:50-10:40 -- Lamanaco, Day 4 -- 9:50-10:40 -- Lallucci, Day 5 -- 9:50-10:40 -- Regelsberger, Day 5 -- 2:15-3:05 -- Tandoi
\item \textbf{Main Lesson Family:} Block Coding with Sphero
\item \textbf{Primary Evidence Base:} Transcript-derived notes plus source transcripts from this teaching-placement workspace.
\item \textbf{Primary External Source Type:} Official curriculum/provider materials.
\end{itemize}

\section*{What We Are Teaching}
\begin{itemize}[leftmargin=*]
\item pairing and aiming a robot
\item sequencing movement blocks into a route
\item using loops or repeated structures to shorten code
\item debugging distance, direction, and timing
\end{itemize}

\section*{CS Standards Alignment}
\begin{itemize}[leftmargin=*]
\item \textbf{1B-AP-10}: Create programs that include sequences, events, loops, and conditionals. Why it fits here: Students build working movement programs with repeated structures.
\item \textbf{1B-AP-11}: Decompose problems into smaller, manageable subproblems. Why it fits here: Students break the route into single moves before combining them.
\item \textbf{1B-AP-15}: Test and debug a program or algorithm to ensure it runs as intended. Why it fits here: Students repeatedly test, adjust, and rerun their robot code.
\end{itemize}

\section*{NYS / Local Curriculum Alignment}
\begin{itemize}[leftmargin=*]
\item Aligned to the local Grade 3 IGNITE robotics sequence in the workbook and schedule.
\item Supports New York-facing computational thinking and collaborative problem solving through route planning, testing, and revision.
\item This is a core CS packet because the content centers on algorithm design, repeated patterns, and debugging.
\end{itemize}

\section*{Lesson Content and Concepts}
\begin{itemize}[leftmargin=*]
\item sequence
\item loop
\item debugging
\item robot control
\item decomposition
\end{itemize}

\section*{Evidence / Proof of Instruction}
\begin{itemize}[leftmargin=*]
\item \textbf{Transcript-derived note:} 260227 Grade 3 Lamanaco Sphero Loops [../../transcript-derived-lessons/260227-grade-3-lamanaco-sphero-loops.md] The lesson note explicitly names pairing, loops, square movement, aiming, and safety.
\item \textbf{Transcript-derived note:} 260227 Grade 3 Lamanaco Block Coding [../../transcript-derived-lessons/260227-grade-3-lamanaco-block-coding.md] The additional note shows the same lesson family emphasized core block coding rather than one isolated challenge.
\item \textbf{Transcript-derived note:} 260306 Grade 3 Lamanaco Square Maze [../../transcript-derived-lessons/260306-grade-3-lamanaco-square-maze.md] Later transcripts show students applying the same logic to more precise route and maze navigation tasks.
\item \textbf{Source transcript:} 260227 Grade 3 Lamanaco Sphero Loops [../../transcripts/260227-grade-3-lamanaco-sphero-loops.md] The source transcript records teacher language around loops, pairing, and safety while students work in robot pairs.
\item \textbf{Likely student proof:} Working square paths, looped code blocks, debug revisions, and safe partner management are direct evidence of instruction.
\end{itemize}

\section*{Assessment Evidence}
\begin{itemize}[leftmargin=*]
\item Student pairs the correct robot and keeps it under control.
\item Student writes a movement sequence that completes the task or gets close enough to debug.
\item Student can explain one fix they made after testing.
\end{itemize}

\section*{UDL and Inclusion Supports}
\subsection*{Multiple Representation}
\begin{itemize}[leftmargin=*]
\item Model the route on the floor and on screen before coding.
\item Use arrows, distance markers, and turn visuals.
\item Show one buggy program and one corrected version side by side.
\end{itemize}

\subsection*{Multiple Action / Expression}
\begin{itemize}[leftmargin=*]
\item Allow coding by dragging blocks, pointing to a route card, or verbally sequencing for a partner.
\item Use clearly assigned pair roles: coder, launcher, observer, recorder.
\item Give students a printed planning strip before the digital build.
\end{itemize}

\subsection*{Multiple Engagement}
\begin{itemize}[leftmargin=*]
\item Use robot motion as immediate feedback for effort and revision.
\item Keep challenges short enough for multiple retry cycles.
\item Frame debugging as the normal path to success.
\end{itemize}

\section*{How We Are Already Making It Inclusive}
\begin{itemize}[leftmargin=*]
\item Transcript evidence shows teacher modeling before student handling of the robots.
\item Partner structures were already present, which supports participation for students who need shared problem solving.
\item The lesson used physical movement and visible outcomes, which lowers abstraction load.
\end{itemize}

\section*{How to Improve Inclusion Next Time}
\begin{itemize}[leftmargin=*]
\item Add route mats with start/end icons and numbered steps.
\item Create a quiet-observer role for students who need to participate without launching every round.
\item Build a common error wall with examples like overshoot, wrong angle, and unpaired robot.
\item Store route photos and code screenshots in \texttt{artifacts/} as proof.
\end{itemize}

\section*{Materials and Setup}
\begin{itemize}[leftmargin=*]
\item Sphero robots and chargers
\item Student devices
\item Floor route or taped path
\item Planning strips or dry-erase route cards
\end{itemize}

\section*{Teacher Moves / Script / Routines}
\begin{itemize}[leftmargin=*]
\item Model pairing and aiming before any free handling begins.
\item Have students predict the route before launch.
\item Stop midway to compare a long program and a looped version.
\item Close with a debug story: what failed and what changed.
\end{itemize}

\section*{Common Errors and Debugging}
\begin{itemize}[leftmargin=*]
\item Students may skip pairing steps and blame the code.
\item Distance and angle errors can look like conceptual failure when they are calibration issues.
\item Partner imbalance can leave one student passive without role structure.
\end{itemize}

\section*{Extensions and Simplifications}
\begin{itemize}[leftmargin=*]
\item Extension: compare two possible solutions and justify which is more efficient.
\item Simplification: reduce the task to one turn and one forward move repeated.
\item Extension: add an event or sensor-based trigger after the movement pattern works.
\end{itemize}

\section*{Related Local Materials}
\begin{itemize}[leftmargin=*]
\item Block Coding with Bees materials [../../lessons/23-Block-Coding-with-Bees/23-Block-Coding-with-Bees-Notes.md]
\end{itemize}

\section*{Artifacts Checklist}
\begin{itemize}[leftmargin=*]
\item Compiled lesson PDF in \texttt{artifacts/}
\item At least one screenshot, student example, or setup photo
\item One short assessment or observation record
\item Updated standards/source bibliography once the \texttt{.bib} file is revised
\item Any teacher reflection or revision notes after reteaching
\end{itemize}

\section*{Selected Official Sources}
Selected official sources that informed this packet are listed below.
ocite{*}

\bibliographystyle{plain}
\bibliography{block-coding-with-sphero}
\end{document}
